\section{Líneas de trabajo futuras}

Las limitaciones descritas en la sección anterior no son todas del mismo tipo:
unas se resuelven con trabajo acotado y otras exigirían rediseñar una parte del
sistema. Se recogen aquí ordenadas por esa distancia, indicando en cada caso qué
limitación concreta vendrían a resolver.

\subsection{Mejoras sobre lo ya construido}

\begin{itemize}

    \item \textbf{Reescritura de la pregunta con el contexto.} Es la respuesta
    al fallo de resolución de referencias entre turnos. Antes de llamar al
    modelo principal, una llamada corta reformularía ``¿y el año pasado?'' como
    una pregunta completa y autónoma a partir de los últimos turnos. Tiene un
    coste claro —una llamada más por turno— que habría que medir antes de
    adoptarlo, ya que la latencia es una prioridad del proyecto.

    \item \textbf{Cambiar a un modelo más capaz si el equipo lo permite.} El
    modelo es una variable de entorno (\texttt{MODELO}, con \texttt{qwen3:8b}
    por defecto), así que sustituirlo no supone tocar el código. Lo que decide
    es el hardware: un modelo mayor necesita más memoria en la tarjeta gráfica y
    genera más despacio, y en este proyecto la latencia es una prioridad. Es,
    además, la otra vía para atacar el fallo de referencias entre turnos, que es
    una limitación de razonamiento propia de un modelo de 8.000 millones de
    parámetros más que un problema de diseño. El banco de pruebas es justamente
    la herramienta para decidirlo: permite comparar precisión y latencia entre
    dos modelos sobre las mismas 44 preguntas y comprobar si la mejora compensa
    el tiempo de respuesta añadido.

    \item \textbf{Previsión de meses distintos del actual.} La herramienta
    \texttt{proyectar\_gasto} calcula siempre sobre el mes en curso. Añadirle un
    parámetro de mes objetivo, y que rechace explícitamente los meses para los
    que no hay base suficiente, evitaría que el modelo presente la cifra del mes
    actual como si fuera la del siguiente.

    \item \textbf{Servir las dependencias del navegador en local.} Descargar
    Vega-Lite, las tipografías y los iconos y servirlos desde
    \texttt{frontend/} haría que la aplicación funcionase sin conexión y
    eliminaría la única dependencia externa que queda.

    \item \textbf{Registro de auditoría.} Guardar en una tabla cada consulta
    SQL generada, cada herramienta invocada y cada operación ejecutada, con su
    resultado y su marca de tiempo. Es poco trabajo y encaja con el enfoque de
    IA responsable: permite reconstruir después por qué el asistente respondió
    lo que respondió.

    \item \textbf{Ampliar el banco de pruebas y automatizarlo.} Las preguntas
    actuales cubren los tipos del enunciado, pero el banco se ejecuta a mano.
    Lanzarlo de forma automática ante cada cambio del \textit{prompt} detectaría
    las regresiones en el momento en que se introducen, que es justo el problema
    que motivó construirlo.

\end{itemize}

\subsection{Mejoras que requieren rediseño}

\begin{itemize}

    \item \textbf{Autenticación real y gestión de sesión.} Sustituir la pantalla
    de acceso por una validación en el servidor, proteger el punto de entrada
    \texttt{/ws} con un testigo de sesión, almacenar el PIN cifrado y limitar el
    número de intentos por origen. Es el paso obligado para que la aplicación
    deje de ser una demostración.

    \item \textbf{Soporte multiusuario.} Implica añadir la noción de cuenta a
    las tablas y a todas las consultas, asociar cada conexión a un titular y
    sustituir SQLite por un gestor que admita varios escritores. El patrón de
    validación, PIN y transacción se mantendría; lo que cambia es de dónde sale
    la cuenta sobre la que se opera.

    \item \textbf{Escalado del modelo.} Con un único modelo en una tarjeta
    gráfica, las peticiones se atienden una detrás de otra. Atender a varios
    usuarios a la vez exigiría una cola explícita y varias réplicas del modelo,
    además de decidir qué ocurre cuando la espera supera lo tolerable.

    \item \textbf{Ampliar el catálogo de operaciones.} Transferencias,
    domiciliaciones o gestión de tarjetas encajarían en el mismo esquema que el
    Bizum: extracción de la intención, validación en el \emph{backend},
    confirmación con PIN y transacción. Cada nueva operación añade, eso sí,
    texto al \textit{system prompt}, y ya se midió que el tamaño del
    \textit{prompt} degrada las reglas existentes.

    \item \textbf{Persistencia de la conversación.} Guardar el historial fuera
    del proceso permitiría retomar una conversación tras una desconexión y
    abriría la puerta a que el asistente recuerde preferencias entre sesiones.

\end{itemize}

De todas ellas, las dos que más cambiarían la experiencia son la reescritura de
la pregunta con el contexto, porque elimina el único fallo sistemático que la
evaluación detecta, y la autenticación real, porque es la diferencia entre una
demostración y algo que pueda tocar dinero de verdad.
